\documentclass[11pt,reqno]{amsart}
\usepackage[T1]{fontenc}
\usepackage{lmodern,microtype}
\usepackage[letterpaper,left=1.12in,right=1.12in,top=1.0in,bottom=1.0in,headheight=14pt]{geometry}
\usepackage{amsmath,amssymb,amsthm,mathtools,mathrsfs}
\usepackage{booktabs,tabularx,array}
\usepackage{xurl}
\usepackage[hidelinks,unicode]{hyperref}
\hypersetup{pdftitle={Elementary positivity of weighted chromatic functions and dual Dyck symmetric functions},
 pdfsubject={Weighted chromatic elementary positivity and step-tau Dyck factorizations},
 pdfkeywords={Dyck symmetric functions, elementary positivity, natural unit interval graphs, multicolorings, common histories}}
\newcommand{\NN}{\mathbb Z_{\geq0}}
\newcommand{\ZZ}{\mathbb Z}
\newcommand{\QQ}{\mathbb Q}
\newcommand{\xx}{\mathbf x}
\newcommand{\zz}{\mathbf z}
\newcommand{\mm}{\mathbf m}
\newcommand{\Sym}{\Lambda}
\newcommand{\QSym}{\operatorname{QSym}}
\newcommand{\DS}{\operatorname{DS}}
\newcommand{\DD}{\mathcal D}
\newcommand{\EE}{\mathscr E}
\newcommand{\CC}{\mathscr C}
\newcommand{\HH}{\mathcal H}
\newcommand{\WW}{\mathcal W}
\newcommand{\asc}{\operatorname{asc}}
\newcommand{\inv}{\operatorname{inv}}
\newcommand{\dinv}{\operatorname{dinv}}
\newcommand{\di}{\operatorname{di}}
\newcommand{\Des}{\operatorname{Des}}
\newcommand{\SYT}{\operatorname{SYT}}
\newcommand{\csf}{\mathsf{csf}}
\newcommand{\supp}{\operatorname{supp}}
\newcommand{\wt}{\operatorname{wt}}
\newcommand{\Frob}{\operatorname{Frob}}
\newcommand{\Ind}{\operatorname{Ind}}
\newcommand{\one}{\mathbf 1}
\newcommand{\sourcefile}[2]{\href{https://github.com/openai/math/blob/fd4aeeb/lean/OAI/Combinatorics/Chromatic/#1}{\nolinkurl{#2}}}
\newcommand{\decl}[1]{\nolinkurl{#1}}
\newtheorem{maintheorem}{Theorem}
\renewcommand{\themaintheorem}{\Alph{maintheorem}}
\newtheorem{theorem}{Theorem}[section]
\newtheorem{proposition}[theorem]{Proposition}
\newtheorem{lemma}[theorem]{Lemma}
\newtheorem{corollary}[theorem]{Corollary}
\theoremstyle{definition}
\newtheorem{definition}[theorem]{Definition}
\newtheorem{example}[theorem]{Example}
\theoremstyle{remark}
\newtheorem{remark}[theorem]{Remark}
\numberwithin{equation}{section}
\setlength{\emergencystretch}{1.5em}
\setlength{\tabcolsep}{5pt}
\renewcommand{\arraystretch}{1.12}
\title[Weighted chromatic and dual Dyck positivity]{Elementary positivity of weighted chromatic functions and dual Dyck symmetric functions}
\author{Graham Hawkes}
\date{October 10, 2026}
\keywords{Dyck symmetric functions, elementary positivity, chromatic quasisymmetric functions, natural unit interval graphs, multicolorings}

\begin{document}
\begin{abstract}
We prove coefficientwise elementary positivity for the dual step-$\tau$
Dyck factorization series of every finite multiset of integers and every
integer $\tau\geq0$.  The corresponding affine series are coefficientwise
positive in the complete homogeneous basis.  With an explicit change of
statistic, the step-one specialization proves elementary positivity of
Hawkes's dual Dyck symmetric functions and homogeneous positivity of his
affine Dyck symmetric functions.  The main algebraic result is a weighted
extension of chromatic elementary positivity: for a natural unit interval
graph, a vertex of weight $m$ receives a set of $m$ colors, and the resulting
inversion-refined chromatic function is elementary-positive over
$\mathbb Z_{\geq0}[u]$.  We use OpenAI's positive common-history theorem as
an external input.  The extension replaces a squarefree vertex target by
an arbitrary multiplicity target and proves the associated coefficient
identity directly.  A squarefree color coefficient then forces each
positive Laurent history weight to have finite support in the correct
grading.  This avoids division by quantum factorials.  We also obtain
palindromicity, a history formula for the elementary coefficients, and
positivity for generalized Hessenberg chromatic functions.
\end{abstract}
\maketitle

\section{Introduction and main results}\label{sec:intro}

The dual and affine Dyck symmetric functions introduced by
Hawkes~\cite{Hawkes2026} arise from factorizations of a multiset into
integer sequences satisfying complementary local inequalities.  His
insertion correspondence proves Schur positivity of the dual functions;
an involution gives the conjugate-shape expansion for the affine
functions~\cite[Theorem~3.18 and Corollaries~3.19--3.20]{Hawkes2026}.
The present paper strengthens these conclusions to positivity in the
elementary and complete homogeneous bases, respectively.  We work
throughout with an integer step $\tau\geq0$, obtaining the original
functions as a step-one corollary with their original statistic.

A dual factor has strictly increasing entries whose successive gaps
exceed $\tau$.  Its set of values is therefore an independent set in the
graph joining distinct values at distance at most $\tau$.  When a value
occurs repeatedly in the multiset, its occurrences lie in distinct
factors.  Thus the appropriate graph model is a \emph{multicoloring}:
a vertex of weight $m$ receives a set of $m$ colors.  This weighted
chromatic construction is already present in Hwang~\cite{Hwang2024}
and Kiem--Lee~\cite[Definitions~4.1--4.3]{KiemLee2024}.
Numbered references to Hawkes and Kiem--Lee use the respective arXiv
versions listed in the bibliography.  It is not the
alternative vertex-weighted construction that assigns one color, raised
to a vertex-dependent power, to each vertex.

The unweighted chromatic quasisymmetric functions of natural unit
interval graphs were introduced by Shareshian and
Wachs~\cite{ShareshianWachs2016}.  Our proof uses the common-history
machinery underlying OpenAI's elementary-positivity
argument~\cite{OpenAI2026,OpenAILean2026}.  The intermediate theorem needed
here is more flexible than the final unweighted statement: it transports
a general root-cone coefficient to a positive combination of ordinary
elementary coefficients.  We exploit this flexibility at a target with
arbitrary vertex multiplicities.

\subsection{The weighted theorem}
A \emph{natural unit interval graph} on $[\ell]=\{1,\ldots,\ell\}$ is a
graph specified by a weakly increasing function
$h:[\ell]\to[\ell]$ with $h(i)\geq i$, where, for $i<j$,
\[
 ij\in E(G)\quad\Longleftrightarrow\quad j\leq h(i).
\]
Fix positive integer weights $\mm=(m_1,\ldots,m_\ell)$ and write
\begin{equation}\label{eq:NA}
 N=\sum_i m_i,\qquad
 A=A(G,\mm)=\sum_{\substack{i<j\\ij\in E(G)}}m_i m_j.
\end{equation}
A proper set coloring $\gamma$ assigns a set $\gamma_i\subset\mathbb Z_{>0}$
of cardinality $m_i$ to vertex $i$, with
$\gamma_i\cap\gamma_j=\varnothing$ whenever $ij\in E(G)$.  Put
\begin{align}
 \inv_G(\gamma)
 &=\sum_{\substack{i<j\\ij\in E(G)}}
   \#\{(a,b)\in\gamma_i\times\gamma_j:a>b\},\label{eq:weighted-inv}\\
 C^-_{G,\mm}(\xx;u)
 &=\sum_{\gamma\text{ proper}}u^{\inv_G(\gamma)}
       \prod_i\prod_{a\in\gamma_i}x_a.\label{eq:Cminus}
\end{align}
The ascent version $C^+_{G,\mm}$ uses $a<b$ in
\eqref{eq:weighted-inv}.  These are homogeneous quasisymmetric series of
degree $N$ with polynomial coefficients in $u$.

\begin{maintheorem}[Weighted chromatic elementary positivity]\label{thm:A}
For every natural unit interval graph $G$ and every positive integer
weight vector $\mm$, there are polynomials
$B_{G,\mm,\lambda}(u)\in\mathbb Z_{\geq0}[u]$, indexed by $\lambda\vdash N$,
such that
\begin{equation}\label{eq:thmA}
 C^+_{G,\mm}(\xx;u)=C^-_{G,\mm}(\xx;u)
   =\sum_{\lambda\vdash N}B_{G,\mm,\lambda}(u)e_\lambda(\xx).
\end{equation}
Every coefficient polynomial has degree at most $A$ and satisfies
\begin{equation}\label{eq:thmA-pal}
 B_{G,\mm,\lambda}(u)=u^A B_{G,\mm,\lambda}(u^{-1}).
\end{equation}
\end{maintheorem}

In particular, symmetry is part of the conclusion, not an input to our
proof.  For $\mm=(1,\ldots,1)$, this recovers the unweighted result.
The general weights will be essential for multisets with repeated values.

\subsection{The Dyck theorem}
Write a finite multiset of integers as
\[
 S=c_1^{m_1}\cdots c_\ell^{m_\ell},\qquad
 c_1<\cdots<c_\ell,\qquad |S|=N.
\]
For a fixed integer $\tau\geq0$, define
\begin{align}
 A_\tau(S)&=\sum_{\substack{i<j\\c_j-c_i\leq\tau}}m_i m_j,
       \label{eq:Atau}\\
 \beta_\tau(S)&=\tau\sum_i\binom{m_i}{2}
  +\sum_{\substack{i<j\\c_j-c_i\leq\tau}}
       (\tau-c_j+c_i)m_i m_j.\label{eq:beta}
\end{align}
The dual and affine factorization series defined in
Section~\ref{sec:definitions} are denoted by
$\DD^*_{S,\tau}(\xx;u)$ and $\DD_{S,\tau}(\xx;u)$.  Their coefficient of
$u^d$ counts factorizations whose concatenation has statistic
$\dinv_\tau=d$.

\begin{maintheorem}[Step-$\tau$ Dyck positivity]\label{thm:B}
For every finite multiset $S$ of integers and every integer $\tau\geq0$,
there are nonnegative integers $a_{\lambda,k}(S,\tau)$ such that
\begin{align}
 \DD^*_{S,\tau}(\xx;u)
  &=\sum_{\lambda\vdash N}\sum_{k=0}^{A_\tau(S)}
       a_{\lambda,k}(S,\tau)u^{\beta_\tau(S)+k}e_\lambda(\xx),
       \label{eq:thmB-dual}\\
 \DD_{S,\tau}(\xx;u)
  &=\sum_{\lambda\vdash N}\sum_{k=0}^{A_\tau(S)}
       a_{\lambda,k}(S,\tau)u^{\beta_\tau(S)+k}h_\lambda(\xx).
       \label{eq:thmB-affine}
\end{align}
Moreover,
\begin{equation}\label{eq:thmB-pal}
 a_{\lambda,k}(S,\tau)
   =a_{\lambda,A_\tau(S)-k}(S,\tau).
\end{equation}
Thus each fixed-statistic dual function is $e$-positive and each
fixed-statistic affine function is $h$-positive.  Both are zero outside
$\beta_\tau(S)\leq d\leq\beta_\tau(S)+A_\tau(S)$.
\end{maintheorem}

For $S=\varnothing$, the assertion means that both series are $1$,
with $a_{\varnothing,0}=1$.  The nonempty case occupies the proof.
At step one, the factor conditions agree with Hawkes's definitions, but
our statistic includes the constant
$\kappa(S)=\sum_i\binom{m_i}{2}$.  After this shift is removed,
Theorem~\ref{thm:B} gives the following precise classical consequence.

\begin{corollary}\label{cor:classical-intro}
For every finite multiset $S$ of integers and every $d\geq0$, the functions
$\DS^*(S,d;\xx)$ and $\DS(S,d;\xx)$ of
\cite[Definition~2.5]{Hawkes2026} are respectively $e$-positive and
$h$-positive, with nonnegative integer coefficients.
\end{corollary}

\subsection{Why the multiplicity extension is needed}
Replace each vertex $i$ of $G$ by an ordered clique of $m_i$ vertices,
and replace each original edge by all edges between the corresponding
cliques.  For the resulting graph $\widetilde G$, the standard blowup
identity is
\begin{equation}\label{eq:blowup-intro}
 X_{\widetilde G}(\xx;u)
    =\left(\prod_i[m_i]_u!\right)C^+_{G,\mm}(\xx;u),
 \qquad [m]_u!=\prod_{j=1}^m(1+u+\cdots+u^{j-1}).
\end{equation}
See \cite{Hwang2024} and \cite[Remark~4.4 and (4.3)]{KiemLee2024};
a direct proof is included below.  Elementary positivity of the left
side does not imply coefficientwise positivity after division by the
factor on the right.  Even
\[
 (1+u)(1-u+u^2)=1+u^3
\]
shows why such cancellation would be invalid.

Instead, we extract the coefficient of a lattice monomial with
multiplicity vector $\mm$ from an independent-set product in a quantum
torus.  The common-history theorem applies to this coefficient before
any blowup factorial is introduced.  Its weights are positive Laurent
series.  A second extraction, this time of the squarefree \emph{color}
monomial $x_1\cdots x_N$, forces their support to lie in the finite set
$\{A,A-2,\ldots,-A\}$.  Transposition of zero-one matrices then supplies
the elementary expansion over $\mathbb Z_{\geq0}[u]$.

Section~\ref{sec:definitions} introduces the step-$\tau$ factorizations
and their grading.  Section~\ref{sec:colorings} proves the static-statistic
identity and the weighted coloring correspondence.  Section~\ref{sec:input} states the precise imported input, and
Section~\ref{sec:weighted-proof} proves Theorem~\ref{thm:A}.
Section~\ref{sec:dyck-proof} completes the Dyck theorem and its classical
specialization.  Section~\ref{sec:examples} records direct consequences
and examples.  Appendix~\ref{app:sources} gives the source correspondence
at a fixed repository revision, and Appendix~\ref{app:checks} describes
independent exact computations.  The connection to $q,t$-Catalan
polynomials in~\cite{Hawkes2026} motivates the result; no additional
Catalan identity is needed for the proofs below.

\section{\texorpdfstring{Step-$\tau$}{Step-tau} factorizations and grading conventions}\label{sec:definitions}

\subsection{Symmetric functions and coefficientwise positivity}
We work with the ring $\Sym_\ZZ$ of symmetric functions in
$\xx=(x_1,x_2,\ldots)$.  The elementary and complete homogeneous functions
are
\[
 e_r(\xx)=\sum_{a_1<\cdots<a_r}x_{a_1}\cdots x_{a_r},\qquad
 h_r(\xx)=\sum_{a_1\leq\cdots\leq a_r}x_{a_1}\cdots x_{a_r},
\]
with $e_0=h_0=1$.  For a partition $\lambda$, write
$e_\lambda=\prod_i e_{\lambda_i}$ and
$h_\lambda=\prod_i h_{\lambda_i}$.  The involution
$\omega:\Sym_\ZZ\to\Sym_\ZZ$ is characterized by
$\omega(e_r)=h_r$; equivalently, $\omega(s_\lambda)=s_{\lambda'}$.
It acts coefficientwise on $\Sym_\ZZ[u]$ and leaves $u$ fixed.

An element $\sum_\lambda b_\lambda(u)e_\lambda$ is
\emph{coefficientwise $e$-positive} if every
$b_\lambda(u)\in\mathbb Z_{\geq0}[u]$.  This is stronger than positivity
after a numerical specialization of $u$.  The analogous definition
uses $h_\lambda$.  Finite-variable restrictions will be denoted by a
subscript $r$ or by displaying $x_1,\ldots,x_r$.  All our series have
fixed homogeneous degree and are determined by these restrictions.

There are three distinct parameters in the proof.  The integer $N$ is
the homogeneous degree in $\xx$, the variable $u$ records the inversion
or Dyck statistic, and the Laurent variable $v$ is used only in the
quantum-torus argument.  Its relation to $u$ is $u=v^{-2}$, after a
normalizing factor $v^A$.  In particular, $v$ is not a second Catalan
parameter.

\subsection{Words and factorizations}
Let $\WW(S)$ be the set of distinct rearrangements of the multiset $S$.
Equal occurrences are not labeled.  Fix $\tau\in\NN$.

\begin{definition}\label{def:words}
A finite integer word $w=(w_1,\ldots,w_n)$ is an \emph{affine step-$\tau$
word} if $w_{j+1}\leq w_j+\tau$ for every $j<n$.  It is a \emph{dual
step-$\tau$ word} if $w_{j+1}>w_j+\tau$ for every $j<n$.
The empty word and each one-letter word satisfy both conditions.
\end{definition}

No nonnegativity or initial-entry condition is imposed.  These are
conditions on individual factors, not on the multiset itself.  A
factorization of $S$ is a sequence
$\mathcal F=(F_1,F_2,\ldots)$ of finite words, all but finitely many empty,
whose concatenation $w(\mathcal F)=F_1F_2\cdots$ lies in $\WW(S)$.
It is affine or dual when each factor has the corresponding property.
Its monomial is
\begin{equation}\label{eq:factor-monomial}
 \xx^{\mathcal F}=\prod_{a\geq1}x_a^{|F_a|}.
\end{equation}
Empty factors are important: factor indices are the color labels that
will be used later.  Our one-indexed convention differs from the
zero-indexed convention in~\cite{Hawkes2026} only by renaming variables
and factor indices.

\subsection{The statistic}
For integers $a,b$, put
\begin{equation}\label{eq:pair-stat}
 d_\tau(a,b)=
 \begin{cases}
  \max(0,a+\tau-b),&a\leq b,\\
  \max(0,b+1+\tau-a),&a>b.
 \end{cases}
\end{equation}
For a word $w$, define
\begin{equation}\label{eq:dinv}
 \dinv_\tau(w)=\sum_{p<q}d_\tau(w_p,w_q).
\end{equation}
The statistic of a factorization is that of its concatenation; there is
no additional contribution attached to a factor boundary.  Define
\begin{align}
 \DD^*_{S,\tau}(\xx;u)
   &=\sum_{\mathcal F\text{ dual}}u^{\dinv_\tau(w(\mathcal F))}
            \xx^{\mathcal F},\label{eq:dual-series}\\
 \DD_{S,\tau}(\xx;u)
   &=\sum_{\mathcal F\text{ affine}}u^{\dinv_\tau(w(\mathcal F))}
            \xx^{\mathcal F}.\label{eq:affine-series}
\end{align}
Their degree-$d$ coefficients are
\begin{equation}\label{eq:degree-d}
 \DS^{*(\tau)}(S,d;\xx)=[u^d]\DD^*_{S,\tau}(\xx;u),\qquad
 \DS^{(\tau)}(S,d;\xx)=[u^d]\DD_{S,\tau}(\xx;u).
\end{equation}
Both are understood to be zero for $d<0$.  Before symmetry has been
established, we refer to~\eqref{eq:dual-series}--\eqref{eq:affine-series}
as factorization series.

Every coefficient of a fixed monomial is a finite sum: its exponents
fix all factor lengths, and only finitely many rearrangements of $S$
can be concatenations.  An increasing relabeling of the used factor
indices preserves concatenation and statistic, so the series are
quasisymmetric.  Their $u$-degrees are uniformly bounded for fixed $S$
and $\tau$, since there are finitely many words in $\WW(S)$.

\begin{remark}\label{rem:integer-step}
The integrality of the entries and of $\tau$ matters in the grading
calculation below.  We prove the full integer-step family; no assertion
about arbitrary rational slopes is implicit in the terminology.
\end{remark}

\section{The weighted coloring model}\label{sec:colorings}

\subsection{Separating the constant part of the statistic}
For $w\in\WW(S)$, set
\begin{equation}\label{eq:inv-tau}
 \inv_\tau(w)=\#\{(p,q):p<q,\ 0<w_p-w_q\leq\tau\}.
\end{equation}

\begin{lemma}[Static decomposition]\label{lem:static}
For every $w\in\WW(S)$,
\begin{equation}\label{eq:static}
 \dinv_\tau(w)=\beta_\tau(S)+\inv_\tau(w).
\end{equation}
\end{lemma}
\begin{proof}
Consider a pair of occurrences, independently of its order in the word.
If both entries equal $c_i$, its contribution to~\eqref{eq:dinv} is
$\tau$.  Summing over equal pairs gives
$\tau\sum_i\binom{m_i}{2}$.

For two distinct values $c_i<c_j$, let $\delta=c_j-c_i$.  If
$\delta>\tau$, both orders contribute zero: for the decreasing order,
$\tau+1-\delta\leq0$ because $\delta$ is an integer.  If
$1\leq\delta\leq\tau$, then
\[
 d_\tau(c_i,c_j)=\tau-\delta,\qquad
 d_\tau(c_j,c_i)=\tau-\delta+1.
\]
There is an order-independent contribution $\tau-\delta$, together
with an additional $1$ exactly when the larger value precedes the
smaller.  There are $m_i m_j$ pairs with these two values.  Summing over
all value pairs proves~\eqref{eq:static}.
\end{proof}

\begin{corollary}\label{cor:support}
Every factorization of $S$ has statistic in the interval
\[
 \beta_\tau(S)\leq\dinv_\tau\leq\beta_\tau(S)+A_\tau(S).
\]
\end{corollary}
\begin{proof}
Exactly $A_\tau(S)$ pairs of occurrences have distinct values within
distance $\tau$.  Each contributes either $0$ or $1$ to
$\inv_\tau$.
\end{proof}

The baseline is not just an equal-entry correction when $\tau>1$:
nearby distinct values also contribute to~\eqref{eq:beta}.  This
observation will account for the different grading shifts in the
examples of Section~\ref{sec:examples}.

\subsection{The value graph and the coloring correspondence}
Define $G_{S,\tau}$ on the value levels $[\ell]$ by
\begin{equation}\label{eq:value-graph}
 ij\in E(G_{S,\tau})\quad\Longleftrightarrow\quad
 i<j\text{ and }c_j-c_i\leq\tau.
\end{equation}
Vertex $i$ has weight $m_i$.

\begin{lemma}\label{lem:interval}
The graph $G_{S,\tau}$ is a natural unit interval graph, with
Hessenberg function
\begin{equation}\label{eq:value-h}
 h_{S,\tau}(i)=\max\{j\in[\ell]:c_j-c_i\leq\tau\}.
\end{equation}
Its weighted edge count is $A_\tau(S)$.
\end{lemma}
\begin{proof}
The maximum exists because $j=i$ is admissible.  Increasing $i$
increases the threshold $c_i+\tau$, so the maximum is weakly increasing.
It is at least $i$.  For $i<j$, condition~\eqref{eq:value-graph} is
exactly $j\leq h_{S,\tau}(i)$.  The assertion about the weighted edge
count follows from~\eqref{eq:Atau}.
\end{proof}

\begin{proposition}[Dual factorizations as set colorings]\label{prop:bijection}
There is a monomial-preserving bijection between dual step-$\tau$
factorizations of $S$ and proper set colorings of $(G_{S,\tau},\mm)$.
Under this bijection,
\begin{equation}\label{eq:stat-bijection}
 \inv_\tau(w(\mathcal F))=\inv_{G_{S,\tau}}(\gamma).
\end{equation}
Consequently,
\begin{equation}\label{eq:dual-color}
 \DD^*_{S,\tau}(\xx;u)
   =u^{\beta_\tau(S)}C^-_{G_{S,\tau},\mm}(\xx;u).
\end{equation}
\end{proposition}
\begin{proof}
A dual factor is strictly increasing, so it contains at most one copy of
each value.  All pairs of its values have distance greater than $\tau$:
this is immediate for consecutive values, and a nonconsecutive
positive gap is even larger.  Thus its value levels form an independent
set of $G_{S,\tau}$.

Given $\mathcal F=(F_1,F_2,\ldots)$, set
\[
 \gamma_i=\{a\geq1:c_i\text{ occurs in }F_a\}.
\]
There are exactly $m_i$ such factors.  Adjacent value levels cannot
occur in the same factor, so $\gamma$ is proper.  Conversely, from a
proper coloring, form $F_a$ by listing, in increasing order, those
$c_i$ for which $a\in\gamma_i$.  Its value levels are independent, and
hence its successive gaps exceed $\tau$.  These constructions are
inverse.  They also give
$|F_a|=\#\{i:a\in\gamma_i\}$, proving monomial preservation.

Consider $i<j$ with $c_j-c_i\leq\tau$, and choose colors
$a\in\gamma_i$, $b\in\gamma_j$.  Properness gives $a\ne b$.
The occurrence of $c_j$ precedes that of $c_i$ in the concatenation
precisely when $b<a$.  Thus it contributes to~\eqref{eq:inv-tau}
precisely when the pair $(a,b)$ contributes to
$\inv_{G_{S,\tau}}(\gamma)$.  There is no contributing pair within a
dual factor.  This proves~\eqref{eq:stat-bijection}.
Combining it with Lemma~\ref{lem:static} proves~\eqref{eq:dual-color}.
\end{proof}

This identity uses inversions directly.  Neither symmetry nor
palindromicity has been used to reverse the grading.

\subsection{The clique blowup identity}
We record a proof of~\eqref{eq:blowup-intro} to specify its exact
normalization.  Let $M_i=m_1+\cdots+m_i$, with $M_0=0$.  The blowup
$\widetilde G$ has the consecutive blocks
$\{M_{i-1}+1,\ldots,M_i\}$.  It is again a natural unit interval graph:
its Hessenberg function is
\begin{equation}\label{eq:blowup-h}
 \widetilde h(a)=M_{h(i)}
 \quad\text{for }M_{i-1}<a\leq M_i.
\end{equation}

\begin{proposition}\label{prop:blowup}
Let $X_{\widetilde G}$ use the usual ascent statistic on proper
single-color assignments to $\widetilde G$.  Then
\begin{equation}\label{eq:blowup}
 X_{\widetilde G}(\xx;u)
 =\left(\prod_i[m_i]_u!\right)C^+_{G,\mm}(\xx;u).
\end{equation}
The corresponding identity also holds with inversions on both sides.
\end{proposition}
\begin{proof}
A proper coloring of a clique uses distinct colors.  Forgetting their
order within each block produces a proper set coloring of $G$.
For a fixed set coloring, the choices of order in different blocks are
independent.  All contributions from edges between blocks are fixed
and equal the weighted ascent statistic.  Within a block of size $m_i$,
the generating polynomial of ascent pairs over all orders of its
colors is $[m_i]_u!$.  Indeed, the inversion generating polynomial of
permutations has the recurrence $P_m=[m]_uP_{m-1}$, obtained by inserting
the largest entry, and reversing the order exchanges ascent and
inversion pairs.  Multiplying the block contributions gives
\eqref{eq:blowup}.  The inversion version has the same proof.
\end{proof}

The proposition explains the relation to ordinary chromatic functions,
but it is not a positivity argument for the quotient.  We now work with
the weighted target itself.

\section{The positive common-history input}\label{sec:input}

This section specifies the external input from
\cite{OpenAI2026,OpenAILean2026}.  We do not reproduce the wall-crossing
and common-history constructions.  Instead, we state exactly the
properties used below and identify the source declarations in
Appendix~\ref{app:sources}.  In particular, the theorem we import is not
merely the final elementary-positivity theorem for unweighted graphs.

\subsection{Quantum-torus notation}
Let $K=\QQ((v))$, and let $L$ be a finite-rank free abelian group with
an integral alternating bilinear form $\Omega$.  Its quantum torus over
$K$ is the vector space with basis $\{T^a:a\in L\}$ and multiplication
\begin{equation}\label{eq:torus}
 T^aT^b=v^{-\Omega(a,b)}T^{a+b}.
\end{equation}
Only finite linear combinations of basis monomials are needed here;
their coefficients may be Laurent series.  The coefficient of $T^R$
in $F$ is denoted by $[T^R]F$.  Bilinearity gives
\begin{equation}\label{eq:torus-product}
 T^{a_1}\cdots T^{a_r}
   =v^{-\sum_{p<q}\Omega(a_p,a_q)}T^{a_1+\cdots+a_r}.
\end{equation}
The source uses a unit called \decl{LaurentRay.vUnit}.  Its $k$th
power has Laurent exponent $-k$, so that unit is $v^{-1}$ in our
notation.  This explains the sign in~\eqref{eq:torus}; see
Appendix~\ref{app:sources}.

Let $V$ be a finite alphabet and let $\iota:\ZZ^V\hookrightarrow L$
be an injective homomorphism.  If $\epsilon_a$ is its standard coordinate
vector, set $u_a=\iota(\epsilon_a)$.  A subset $I\subseteq V$ is
\emph{independent} if $\Omega(u_a,u_b)=0$ for all $a,b\in I$.
Define
\begin{equation}\label{eq:independent}
 \EE_k=\sum_{\substack{I\subseteq V\text{ independent}\\|I|=k}}
       T^{\sum_{a\in I}u_a},\qquad \EE_0=1.
\end{equation}
In particular, $\EE_k=0$ for $k>|V|$.  These are the
independent-set elements of the source, with the displayed
normalization.  A contributing set occurs with coefficient $1$;
there is no normalization by $[k]_u!$ or by $k!$.

\subsection{Triangular realizations}
The triangular construction in~\cite{OpenAILean2026} supplies a finite
alphabet $V=V_D$, a lattice $L=L_D$, an alternating form $\Omega$, an
injective coordinate map $\iota$, an anchor $a_0\in L$, and a finite
list of root vectors $\rho_j\in L$, $j\in J$.  Write
\begin{equation}\label{eq:root-monoid}
 \CC=\left\{\sum_{j\in J}n_j\rho_j:n_j\in\NN\right\}.
\end{equation}
Thus membership in $\CC$ means a nonnegative \emph{integral}
combination, not just membership in the real cone spanned by the roots.

\begin{proposition}[Triangular graph realization, OpenAI]\label{prop:embedding}
For every natural unit interval graph $G$ on $[\ell]$, one can choose
the triangular data above and an injection $b:[\ell]\hookrightarrow V$
such that
\begin{equation}\label{eq:graph-pairing}
 \Omega(u_{b(i)},u_{b(j)})
   =-\one_{\{ij\in E(G)\}}\qquad(i<j).
\end{equation}
The coordinate map $\iota$ is injective on all of $\ZZ^V$, including
coordinates outside the embedded graph.
\end{proposition}

This is the combination of the source's
\decl{graph_triangular_embedding}, \decl{graphVertex_pair}, and
\decl{includeVertices_injective}.  We always choose the actual
triangular realization supplied by that construction.  The next
theorem is asserted for those data, not for arbitrary lattice data
satisfying only~\eqref{eq:graph-pairing}.

\begin{theorem}[Positive common histories, OpenAI]\label{thm:input}
Fix such a triangular realization.  Let $N\in\NN$ and $R\in L$ satisfy
\begin{equation}\label{eq:root-hyp}
 R-Na_0\in\CC.
\end{equation}
There are a finite set $\HH$, profiles
$\mu(\eta)=(\mu_0(\eta),\ldots,\mu_D(\eta))\in\NN^{D+1}$, and weights
$W_\eta(v)\in\QQ((v))$ such that
\begin{equation}\label{eq:profiles}
 \sum_{j=0}^D\mu_j(\eta)=N,\qquad
 [v^s]W_\eta(v)\in\mathbb Z_{\geq0}\quad(s\in\ZZ).
\end{equation}
For every $r\geq1$ and every $\alpha=(\alpha_1,\ldots,\alpha_r)\in\NN^r$
with $\sum_a\alpha_a=N$,
\begin{equation}\label{eq:history}
 [T^R]\EE_{\alpha_1}\cdots\EE_{\alpha_r}
  =\sum_{\eta\in\HH}W_\eta(v)
       [\zz^{\mu(\eta)}]\prod_{a=1}^r e_{\alpha_a}(z_0,\ldots,z_D).
\end{equation}
The same set, profiles, and weights work simultaneously for all $r$
and $\alpha$.
\end{theorem}

Here $[v^s]$ is the ordinary Laurent coefficient, allowed at any
integer exponent.  Theorem~\ref{thm:input} initially does \emph{not}
assert that the weights are Laurent polynomials, nor that their
exponents have a prescribed parity.  Both facts will be proved for
our targets in Section~\ref{sec:weighted-proof}.

\begin{proof}[Identification with the cited theorem]
The source's \decl{triangular_common_history} is quantified over every
homogeneous elementary expression of degree $N$.  Take the expression
whose ordered sequence of elementary indices is
$(\alpha_1,\ldots,\alpha_r)$.  Its torus evaluation is the product on
the left of~\eqref{eq:history}, and its ordinary evaluation is the
product on the right.  This is the same specialization used in the
source's unweighted \decl{graph_common_history} proof.

For a representation $R-Na_0=\sum_j n_j\rho_j$, the source's root
degree is $d=\sum_j n_j$.  Thus~\eqref{eq:root-hyp} provides exactly
its hypothesis \decl{HasRootDegree}, with no squarefreeness condition
on $R$.  The finite index set in~\eqref{eq:history} consists of the
anchored leaves of the source's common history.  Their profiles have
total degree $N$; in fact they are weakly decreasing, a property we
will not need.  Finiteness follows from \decl{leafFintype}, and
\decl{leafWeight_positive} gives positivity of every weight.  Finally,
the source defines \decl{LaurentPositive.Positive} by requiring each
coefficient to equal a natural number.  It therefore gives exactly the
integer positivity in~\eqref{eq:profiles}.  The common-history theorem
chooses its history before quantifying over elementary expressions,
which gives the required uniformity in $r$ and $\alpha$.
\end{proof}

The product form~\eqref{eq:history} is enough for the rest of the
paper.  No assertion about commutativity of arbitrary torus elements
is needed: the products on the left retain the displayed order.

\section{Multiplicity targets and weighted elementary positivity}\label{sec:weighted-proof}

For $N=0$, both weighted series are $1$, so Theorem~\ref{thm:A}
is immediate.  Fix $G$ and $\mm$ with $N>0$, and choose
the realization from Proposition~\ref{prop:embedding}.  All history
data in this section are attached to this one realization and one
target.

\subsection{The root-cone hypothesis for multiplicities}
Set
\begin{equation}\label{eq:target}
 R_{\mm}=\sum_{i=1}^{\ell}m_i u_{b(i)}.
\end{equation}

\begin{lemma}[Admissibility of the multiplicity target]\label{lem:cone}
The target $R_{\mm}$ satisfies
\begin{equation}\label{eq:cone}
 R_{\mm}-Na_0\in\CC.
\end{equation}
\end{lemma}
\begin{proof}
Apply the source's \decl{squarefree_rootDegree} to the injection from
a one-element set to the triangular vertex $b(i)$.  Its conclusion is
\[
 u_{b(i)}-a_0=\sum_{j\in J}n_{ij}\rho_j
 \quad\text{with }n_{ij}\in\NN.
\]
This is the singleton case of the cited result; it imposes no
restriction on how often that singleton may be used in a subsequent
sum.  Therefore
\[
 R_{\mm}-Na_0
  =\sum_i m_i(u_{b(i)}-a_0)
  =\sum_{j\in J}\left(\sum_i m_i n_{ij}\right)\rho_j\in\CC.
\]
More explicitly, if the singleton root degree is $d_i=\sum_j n_{ij}$,
then the displayed expression has root degree $\sum_i m_i d_i$.
This verifies the precise hypothesis of Theorem~\ref{thm:input}.
\end{proof}

The use of a squarefree \emph{singleton} result in this lemma should
not be confused with a squarefreeness restriction on the final target.
The coefficients $m_i$ in~\eqref{eq:target} are arbitrary positive
integers.

\subsection{The weighted coefficient identity}
For $r\geq1$, let $C^-_r$ be the restriction of $C^-_{G,\mm}$ to
colors in $[r]$.  Recall the weighted edge count $A$ from~\eqref{eq:NA}.

\begin{lemma}[Weighted coefficient extraction]\label{lem:kernel}
For $\alpha\in\NN^r$ with $\sum_a\alpha_a=N$,
\begin{equation}\label{eq:kernel}
 [T^{R_{\mm}}]\EE_{\alpha_1}\cdots\EE_{\alpha_r}
     =v^A[\xx^\alpha]C^-_r(\xx;v^{-2}).
\end{equation}
\end{lemma}
\begin{proof}
Expanding~\eqref{eq:independent} chooses an independent subset
$I_a\subseteq V$ of cardinality $\alpha_a$ for each color $a\in[r]$.
Its total lattice exponent is
\[
 \sum_{a=1}^r\sum_{z\in I_a}u_z.
\]
Let $n_z=\#\{a:z\in I_a\}$.  Since $u_z=\iota(\epsilon_z)$ and
$\iota$ is injective, this exponent equals $R_{\mm}$ if and only if
\begin{equation}\label{eq:coordinate-equality}
 n_z=\begin{cases}m_i,&z=b(i),\\0,&z\notin b([\ell]).\end{cases}
\end{equation}
All $n_z$ are nonnegative, so no occurrence of an outside vertex can
cancel another coordinate.

For a contributing tuple, define
\begin{equation}\label{eq:kernel-bijection}
 \gamma_i=\{a\in[r]:b(i)\in I_a\}.
\end{equation}
Equation~\eqref{eq:coordinate-equality} gives $|\gamma_i|=m_i$.
By~\eqref{eq:graph-pairing}, independence of $I_a$ is equivalent to
properness for vertices receiving color $a$.  Thus $\gamma$ is a proper
set coloring, and $|I_a|=\alpha_a$ says that its monomial is
$\xx^\alpha$.  Conversely, such a coloring gives
$I_a=\{b(i):a\in\gamma_i\}$.  This is a bijection, including when
some $\alpha_a=0$.  In particular, it does not order or label the
$m_i$ occurrences at vertex $i$; no factorial factor appears.

It remains to compute the torus weight.  Write
$v_a^{\mathrm{lat}}=\sum_{z\in I_a}u_z$ for the lattice exponent of the
$a$th factor, to distinguish it from the Laurent variable $v$.
By~\eqref{eq:torus-product}, the coefficient contributed by the tuple
is
\begin{equation}\label{eq:kernel-weight}
 v^{-\sum_{a<b}\Omega(v_a^{\mathrm{lat}},v_b^{\mathrm{lat}})}.
\end{equation}
Consider an edge $i<j$ and a pair of colors
$a\in\gamma_i$, $b\in\gamma_j$.  Properness excludes $a=b$.
If $a<b$, the ordered pairing is
$\Omega(u_{b(i)},u_{b(j)})=-1$, contributing $+1$ to the exponent
in~\eqref{eq:kernel-weight}.  If $a>b$, the earlier factor contributes
$u_{b(j)}$ and the later one contributes $u_{b(i)}$; the pairing is
$+1$, contributing $-1$ to that exponent.  Hence the edge contributes
\[
 m_i m_j-2\#\{(a,b)\in\gamma_i\times\gamma_j:a>b\}.
\]
A nonedge contributes zero by~\eqref{eq:graph-pairing}.  Two occurrences
of the same vertex also contribute zero because $\Omega$ is
alternating.  Summing over all edges gives
$A-2\inv_G(\gamma)$.  Thus the term associated to $\gamma$ has weight
$v^A(v^{-2})^{\inv_G(\gamma)}$, proving~\eqref{eq:kernel}.
\end{proof}

The lemma is the new interpretation of coefficient extraction required
for multiplicities.  The common-history theorem itself is unchanged.

\subsection{From common histories to elementary products}
Apply Theorem~\ref{thm:input} at $R_{\mm}$, using
Lemma~\ref{lem:cone}.  Fix the resulting $\HH$, $\mu(\eta)$, and
$W_\eta(v)$.  These do not depend on the number of colors.

\begin{lemma}[Zero-one matrix transposition]\label{lem:transpose}
If $\alpha\in\NN^r$ and $\mu\in\NN^{D+1}$ have the same total size,
then
\begin{equation}\label{eq:transpose}
 [\zz^\mu]\prod_{a=1}^r e_{\alpha_a}(z_0,\ldots,z_D)
 = [\xx^\alpha]\prod_{j=0}^D e_{\mu_j}(x_1,\ldots,x_r).
\end{equation}
\end{lemma}
\begin{proof}
On the left, choose a subset of $\{0,\ldots,D\}$ of size $\alpha_a$
for each factor indexed by $a$.  Its indicator vectors form an
$r\times(D+1)$ zero-one matrix with row sums $\alpha$ and column sums
$\mu$.  On the right, the choices give the transpose matrices, with
column sums $\alpha$ and row sums $\mu$.  Transposition is a bijection.
Zero entries in either profile cause no change to the argument.
\end{proof}

For each history $\eta$, let $\lambda(\eta)\vdash N$ be the partition
obtained by sorting the nonzero entries of $\mu(\eta)$.  Combining
\eqref{eq:history}, \eqref{eq:kernel}, and~\eqref{eq:transpose}, for
every $\alpha$ of total degree $N$, gives
\begin{align*}
 v^A[\xx^\alpha]C^-_r(\xx;v^{-2})
 &=\sum_{\eta\in\HH}W_\eta(v)
          [\zz^{\mu(\eta)}]\prod_{a=1}^r e_{\alpha_a}(\zz)\\
 &=\sum_{\eta\in\HH}W_\eta(v)
          [\xx^\alpha]e_{\lambda(\eta)}(x_1,\ldots,x_r).
\end{align*}
Both sides concern homogeneous polynomials of degree $N$ in the color
variables.  Equality of all coefficients therefore yields
\begin{equation}\label{eq:laurent-expansion}
 v^A C^-_r(\xx;v^{-2})
    =\sum_{\eta\in\HH}W_\eta(v)e_{\lambda(\eta)}(x_1,\ldots,x_r).
\end{equation}
This already gives symmetry over $K$, but we still need an expansion
over $\mathbb Z_{\geq0}[u]$, not merely over positive Laurent series.

\subsection{A positive test coefficient forces finite support}
For a Laurent series $F(v)$, let
$\supp_v(F)=\{s\in\ZZ:[v^s]F\ne0\}$.

\begin{lemma}[Support of a positive sum]\label{lem:positive-support}
Let $F_1,\ldots,F_t\in\QQ((v))$ have nonnegative coefficients, and
let $c_1,\ldots,c_t$ be strictly positive rational numbers.  If
$\sum_i c_iF_i$ is supported on a finite set $E\subseteq\ZZ$, then
$\supp_v(F_i)\subseteq E$ for every $i$.
\end{lemma}
\begin{proof}
For $s\notin E$, coefficient extraction gives
$0=\sum_i c_i[v^s]F_i$.  Each summand is nonnegative.  Since each
$c_i>0$, all the coefficients $[v^s]F_i$ are zero.
\end{proof}

\begin{lemma}[Support and parity of history weights]\label{lem:finite}
For every $\eta\in\HH$,
\begin{equation}\label{eq:weight-support}
 \supp_v(W_\eta)\subseteq E_A:=\{A-2k:0\leq k\leq A\}.
\end{equation}
Consequently there is a polynomial $P_\eta(u)\in\mathbb Z_{\geq0}[u]$
of degree at most $A$ such that
\begin{equation}\label{eq:weight-polynomial}
 W_\eta(v)=v^A P_\eta(v^{-2}).
\end{equation}
\end{lemma}
\begin{proof}
Take $r=N$ in~\eqref{eq:laurent-expansion} and extract the coefficient
of $x_1\cdots x_N$.  Because $\sum_j\mu_j(\eta)=N$, the coefficient
in the elementary product is
\begin{equation}\label{eq:positive-test}
 M_\eta=[x_1\cdots x_N]\prod_{j=0}^D e_{\mu_j(\eta)}(x_1,\ldots,x_N)
       =\frac{N!}{\prod_{j=0}^D\mu_j(\eta)!}>0.
\end{equation}
Indeed, it counts ordered disjoint subsets of the $N$ colors with the
specified cardinalities.  Thus
\begin{equation}\label{eq:test-identity}
 \sum_{\substack{\gamma\text{ proper in }[N]\\\xx^\gamma=x_1\cdots x_N}}
       v^{A-2\inv_G(\gamma)}
   =\sum_{\eta\in\HH}M_\eta W_\eta(v).
\end{equation}
The left side is a finite Laurent polynomial.  Since an edge contributes
between $0$ and $m_i m_j$ inversions,
$0\leq\inv_G(\gamma)\leq A$, and its support lies in $E_A$.
All Laurent coefficients of each $W_\eta$ are nonnegative integers by
Theorem~\ref{thm:input}.  Apply Lemma~\ref{lem:positive-support}, using
the strict positivity of~\eqref{eq:positive-test}, to
obtain~\eqref{eq:weight-support}.

Set
\begin{equation}\label{eq:Peta}
 P_\eta(u)=\sum_{k=0}^A [v^{A-2k}]W_\eta(v)\,u^k.
\end{equation}
Its coefficients are nonnegative integers.  The support statement says
that these are all the coefficients of $W_\eta$, which proves
\eqref{eq:weight-polynomial}.
\end{proof}

The test coefficient in this argument is squarefree in the \emph{color}
variables.  It does not make the lattice target $R_{\mm}$ squarefree:
a vertex of weight $m_i$ still receives $m_i$ distinct colors.  The two
coefficient extractions serve different purposes and are compatible.
The same positive-support principle occurs in the unweighted source
argument; the proof above gives its complete weighted form.

\subsection{Completion of the weighted theorem}
Substitute~\eqref{eq:weight-polynomial} into
\eqref{eq:laurent-expansion} and cancel the invertible monomial $v^A$:
\[
 C^-_r(\xx;v^{-2})
      =\sum_{\eta\in\HH}P_\eta(v^{-2})
                      e_{\lambda(\eta)}(x_1,\ldots,x_r).
\]
The substitution homomorphism $\ZZ[u]\to\QQ((v))$, $u\mapsto v^{-2}$,
is injective: distinct powers of $u$ map to distinct Laurent powers.
Coefficientwise in the color variables, the preceding identity gives
\begin{equation}\label{eq:positive-finite}
 C^-_r(\xx;u)
     =\sum_{\eta\in\HH}P_\eta(u)e_{\lambda(\eta)}(x_1,\ldots,x_r).
\end{equation}
The right side uses the same history polynomials for every $r$.
The restrictions are compatible under $x_r=0$, so passage to symmetric
functions gives
\begin{equation}\label{eq:weighted-expansion}
 C^-_{G,\mm}(\xx;u)
   =\sum_{\lambda\vdash N}B_{G,\mm,\lambda}(u)e_\lambda(\xx),
\end{equation}
where
\begin{equation}\label{eq:history-coeff}
 B_{G,\mm,\lambda}(u)
  =\sum_{\substack{\eta\in\HH\\\lambda(\eta)=\lambda}}
      \sum_{k=0}^A[v^{A-2k}]W_\eta(v)\,u^k
      \in\mathbb Z_{\geq0}[u].
\end{equation}
This proves elementary positivity, integrality, symmetry, and the degree
bound.  The history family is not asserted to be canonical, but the
polynomials $B_{G,\mm,\lambda}$ are canonical by uniqueness of the
elementary expansion.  Formula~\eqref{eq:history-coeff} is a positive
history formula, not a claim of a new tableau or permutation formula
for these coefficients.

It remains to relate the orientations and prove palindromicity.
For a proper set coloring in $[r]$, replace every color $a$ by $r+1-a$.
This is an involution that interchanges ascents and inversions and
reverses the color variables.  Hence
\[
 C^+_r(x_1,\ldots,x_r;u)=C^-_r(x_r,\ldots,x_1;u).
\]
The symmetry just proved makes the right side equal to
$C^-_r(x_1,\ldots,x_r;u)$.  Thus $C^+=C^-$ in the stable limit.
On each edge, equality of colors is prohibited, so
\begin{equation}\label{eq:asc-inv}
 \asc_G(\gamma)+\inv_G(\gamma)=A.
\end{equation}
Termwise, it follows that
\[
 C^+_{G,\mm}(\xx;u)=u^A C^-_{G,\mm}(\xx;u^{-1}).
\]
Comparing elementary coefficients in~\eqref{eq:weighted-expansion}
now proves~\eqref{eq:thmA-pal}.  This completes the proof of
Theorem~\ref{thm:A}.

\begin{remark}\label{rem:no-division}
The only factor cancelled in the positivity argument is the monomial
$v^A$.  No quantum factorial, ordinary factorial, or polynomial with
more than one nonzero coefficient is cancelled from a positive
expansion.  The factorials in~\eqref{eq:positive-test} are merely the
strictly positive scalar test coefficients; they are not used to
divide history weights.
\end{remark}

\section{Dyck positivity and the classical specialization}\label{sec:dyck-proof}

\subsection{The dual series}
Apply Theorem~\ref{thm:A} to the natural unit interval graph
$G_{S,\tau}$ of Lemma~\ref{lem:interval}, with weights
$\mm=(m_1,\ldots,m_\ell)$.  Its weighted edge count is $A_\tau(S)$.
By Proposition~\ref{prop:bijection},
\begin{equation}\label{eq:dual-positive}
 \DD^*_{S,\tau}(\xx;u)
   =u^{\beta_\tau(S)}
       \sum_{\lambda\vdash N}B_{G_{S,\tau},\mm,\lambda}(u)e_\lambda(\xx).
\end{equation}
This proves symmetry and the dual expansion in Theorem~\ref{thm:B}.
Explicitly, we may take
\begin{equation}\label{eq:a-history}
 a_{\lambda,k}(S,\tau)
  =[u^k]B_{G_{S,\tau},\mm,\lambda}(u)
  =\sum_{\substack{\eta\in\HH\\\lambda(\eta)=\lambda}}
       [v^{A_\tau(S)-2k}]W_\eta(v).
\end{equation}
The histories here are chosen for $G_{S,\tau}$ and $R_{\mm}$ as in
Section~\ref{sec:weighted-proof}.  Formula~\eqref{eq:a-history} gives
nonnegative integers and proves the asserted support and
palindromicity of these coefficients.

\subsection{Complementary factor boundaries}
To pass to the affine series, we prove the involution relation for
all integer steps.  This also avoids assuming affine symmetry in
advance.  The argument is the fixed-word complement argument appearing
in the step-one setting in~\cite[Corollary~3.20]{Hawkes2026}.

For $D\subseteq[N-1]$, define the fundamental quasisymmetric function
\begin{equation}\label{eq:fundamental}
 F_{N,D}(\xx)
  =\sum_{\substack{1\leq a_1\leq\cdots\leq a_N\\
                  j\in D\ \Rightarrow\ a_j<a_{j+1}}}
          x_{a_1}\cdots x_{a_N}.
\end{equation}
These functions form a basis of the degree-$N$ quasisymmetric functions.
Let $\psi$ be the degreewise linear involution determined by
\begin{equation}\label{eq:psi}
 \psi(F_{N,D})=F_{N,[N-1]\setminus D}.
\end{equation}
On degree zero put $\psi(1)=1$.

\begin{lemma}\label{lem:psi}
The map $\psi$ preserves symmetric functions and restricts there to
$\omega$.
\end{lemma}
\begin{proof}
For a standard Young tableau $T$ in English convention, let $\Des(T)$
consist of the integers $j$ for which $j+1$ lies in a lower row than $j$.
Standardization of semistandard tableaux gives
\begin{equation}\label{eq:schur-fundamental}
 s_\lambda=\sum_{T\in\SYT(\lambda)}F_{N,\Des(T)}.
\end{equation}
In this standardization, entries of the same value are numbered from
left to right, so equality is allowed exactly away from tableau
descents; this gives~\eqref{eq:schur-fundamental} directly from the
semistandard generating function.  For consecutive entries in a
standard tableau, transposition complements the descent set:
\[
 \Des(T^{\mathsf t})=[N-1]\setminus\Des(T).
\]
Indeed, if $j+1$ lies in a lower row, it cannot lie strictly to the
right of $j$: a cell on a row--column path between the two would have
to carry an integer strictly between $j$ and $j+1$.  If it does not
lie in a lower row, it must lie to the right.  Applying~\eqref{eq:psi}
to~\eqref{eq:schur-fundamental} therefore gives
$\psi(s_\lambda)=s_{\lambda'}=\omega(s_\lambda)$.  Since Schur
functions form a basis of $\Sym_\ZZ$, the assertion follows.
\end{proof}

\begin{proposition}[The affine--dual involution]\label{prop:omega}
For every $S$ and integer $\tau\geq0$,
\begin{equation}\label{eq:omega-dyck}
 \DD_{S,\tau}(\xx;u)=\omega\DD^*_{S,\tau}(\xx;u).
\end{equation}
The identity holds separately in each degree of $u$.
\end{proposition}
\begin{proof}
Fix $w\in\WW(S)$ and write
\begin{align*}
 D^*_\tau(w)&=\{j\in[N-1]:w_{j+1}\leq w_j+\tau\},\\
 D_\tau(w)&=\{j\in[N-1]:w_{j+1}>w_j+\tau\}.
\end{align*}
A factorization with concatenation $w$ is equivalently a weakly
increasing sequence $(a_1,\ldots,a_N)$ of factor labels.  Whenever
$a_j=a_{j+1}$, the pair $(w_j,w_{j+1})$ lies within one factor and
must satisfy its local inequality.  Thus dual factorizations force
a strict increase exactly at the positions in $D^*_\tau(w)$, while
affine factorizations force one exactly at the positions in
$D_\tau(w)$.  Additional strict increases are allowed in either case.
Empty factors correspond to skipped labels, so no further choices
or multiplicities are missing.  The fixed-word generating functions
are consequently $F_{N,D^*_\tau(w)}$ and $F_{N,D_\tau(w)}$.

The two boundary sets are complements.  Since the exponent
$\dinv_\tau(w)$ depends on the concatenated word alone, we obtain
\begin{align*}
 \DD^*_{S,\tau}(\xx;u)
   &=\sum_{w\in\WW(S)}u^{\dinv_\tau(w)}F_{N,D^*_\tau(w)}(\xx),\\
 \DD_{S,\tau}(\xx;u)
   &=\sum_{w\in\WW(S)}u^{\dinv_\tau(w)}F_{N,D_\tau(w)}(\xx)
     =\psi\DD^*_{S,\tau}(\xx;u).
\end{align*}
The dual series is symmetric by~\eqref{eq:dual-positive}, so
Lemma~\ref{lem:psi} replaces $\psi$ by $\omega$.  This proves
\eqref{eq:omega-dyck}, including its coefficientwise version.
For $S=\varnothing$, both sides are $1$.
\end{proof}

Applying $\omega(e_\lambda)=h_\lambda$ to~\eqref{eq:dual-positive}
gives the affine expansion with precisely the same integers
\eqref{eq:a-history}.  This completes the proof of
Theorem~\ref{thm:B}.

\subsection{Hawkes's original functions}
We now identify the step-one specialization with the exact functions
in~\cite{Hawkes2026}.  In Definition~2.1 of that paper, the statistic is
\begin{equation}\label{eq:hawkes-di}
 \di(w)=\#\{p<q:w_p=w_q+1\}.
\end{equation}
Definitions~2.2--2.5 use the affine and dual factor inequalities of
Definition~\ref{def:words} at $\tau=1$.  In particular, equal pairs
are not counted by~\eqref{eq:hawkes-di}.  Put
\begin{equation}\label{eq:kappa}
 \kappa(S)=\sum_i\binom{m_i}{2}.
\end{equation}
Since all entries are integers, $0<w_p-w_q\leq1$ is equivalent to
$w_p=w_q+1$.  Also, the distinct-value terms in $\beta_1(S)$ all
vanish.  Lemma~\ref{lem:static} therefore gives
\begin{equation}\label{eq:classical-shift}
 \beta_1(S)=\kappa(S),\qquad
 \dinv_1(w)=\kappa(S)+\di(w).
\end{equation}
After renaming Hawkes's variable $x_{a-1}$ as our $x_a$, the fixed-degree
identifications are
\begin{align}
 \DS^*_{\mathrm H}(S,d;\xx)
    &=\DS^{*(1)}(S,d+\kappa(S);\xx),\label{eq:classical-dual}\\
 \DS_{\mathrm H}(S,d;\xx)
    &=\DS^{(1)}(S,d+\kappa(S);\xx),\label{eq:classical-affine}
\end{align}
where the subscript $\mathrm H$ denotes the original notation in
\cite[Definition~2.5]{Hawkes2026}, not a different family.

\begin{corollary}[Positivity of the original Dyck functions]\label{cor:classical}
For every finite multiset $S$ and every $d\geq0$,
\begin{equation}\label{eq:classical-expansion}
 \DS^*_{\mathrm H}(S,d;\xx)
    =\sum_{\lambda\vdash N}a_{\lambda,d}(S,1)e_\lambda(\xx),
 \qquad
 \DS_{\mathrm H}(S,d;\xx)
    =\sum_{\lambda\vdash N}a_{\lambda,d}(S,1)h_\lambda(\xx),
\end{equation}
where coefficients outside $0\leq d\leq A_1(S)$ are zero.  All these
coefficients are nonnegative integers, and
$a_{\lambda,d}(S,1)=a_{\lambda,A_1(S)-d}(S,1)$.
\end{corollary}
\begin{proof}
Insert~\eqref{eq:classical-dual}--\eqref{eq:classical-affine} into
Theorem~\ref{thm:B} and use $\beta_1(S)=\kappa(S)$.
\end{proof}

Thus the original generating series equal $u^{-\kappa(S)}\DD^*_{S,1}$
and $u^{-\kappa(S)}\DD_{S,1}$.  These expressions are polynomials in
$u$ by the established lower degree bound.  The removal of this
monomial shift is unrelated to, and does not entail, division by a
quantum factorial.

\section{Consequences and examples}\label{sec:examples}

\subsection{Generalized Hessenberg chromatic functions}
Theorem~\ref{thm:A} also has a direct formulation in the language
of~\cite{KiemLee2024}.  Let
\[
 P=\{M_1<\cdots<M_\ell=N\},\qquad M_0=0,
\]
and let $H:P\to P$ be weakly increasing and extensive.  This is a
generalized Hessenberg function.  Define
\begin{equation}\label{eq:generalized-dictionary}
 m_i=M_i-M_{i-1},\qquad H(M_i)=M_{h(i)}.
\end{equation}
Then $h:[\ell]\to[\ell]$ is weakly increasing and extensive.  Conversely,
any $h$ and positive weights determine $P$ and $H$ in this way.
The weighted graph in~\cite[Definitions~4.1--4.3]{KiemLee2024} joins
$M_i$ to $M_j$ for $M_i<M_j\leq H(M_i)$ and gives $M_i$ the weight
$m_i$.  Its coloring sets, monomial, and ascent count are exactly
those defining $C^+_{G,\mm}$.  Consequently,
\begin{equation}\label{eq:generalized-csf}
 \csf_u(H)=C^+_{G,\mm}(\xx;u).
\end{equation}

\begin{corollary}\label{cor:generalized}
The generalized Hessenberg chromatic function $\csf_u(H)$ is
coefficientwise $e$-positive for every generalized Hessenberg function
$H$.  The equivariant Poincar\'e polynomial $\mathcal F(H)$ in the
normalization of~\cite{KiemLee2024} is coefficientwise $h$-positive.
\end{corollary}
\begin{proof}
The first assertion follows from~\eqref{eq:generalized-csf} and
Theorem~\ref{thm:A}.  The identity
$\csf_u(H)=\omega\mathcal F(H)$ of
\cite[Theorem~4.6]{KiemLee2024} gives the second.
\end{proof}

This answers the positivity question posed in
\cite[Question~4.9]{KiemLee2024}, with the common-history input specified
in Section~\ref{sec:input}.  No geometric result beyond the cited
character identity is needed for this consequence.

\subsection{Boundary cases}
For the empty multiset, both series are $1$.  For $\tau=0$, the value
graph has no edges and both $A_0(S)$ and $\beta_0(S)$ vanish.  Each
vertex chooses its color set independently, giving
\begin{equation}\label{eq:step-zero}
 \DD^*_{S,0}=\prod_i e_{m_i},\qquad
 \DD_{S,0}=\prod_i h_{m_i}.
\end{equation}
For a single value repeated $m$ times, the graph also has no edges,
but the equal-pair shift remains:
\begin{equation}\label{eq:one-value}
 \DD^*_{c^m,\tau}=u^{\tau\binom m2}e_m,\qquad
 \DD_{c^m,\tau}=u^{\tau\binom m2}h_m.
\end{equation}
These formulas in particular distinguish the shift from the
weighted-edge statistic.

At the opposite extreme, if $c_\ell-c_1\leq\tau$, the value graph
is complete.  Every color used in a proper set coloring is then used
once.  Order these $N$ colors increasingly, and record the vertex
receiving each one.  The resulting word has content
$1^{m_1}\cdots\ell^{m_\ell}$, and its ordinary inversion number equals
$\inv_G(\gamma)$.  Hence
\begin{equation}\label{eq:complete}
 C^-_{G,\mm}(\xx;u)
   =\begin{bmatrix}N\\m_1,\ldots,m_\ell\end{bmatrix}_u e_N(\xx),
\end{equation}
where the Gaussian multinomial is defined here by the inversion
generating polynomial of these multiset words.  This directly proves
its nonnegative coefficients.  Its familiar factorial quotient is an
identity for that polynomial, not a positivity inference used in
our proof.

\subsection{A repeated-value example at two steps}
Let
\[
 S=\{0,0,1,2\},\qquad \mm=(2,1,1),\qquad N=4.
\]
At $\tau=1$, the value graph is the path $0$--$1$--$2$, with
$h=(2,3,3)$, $A_1(S)=3$, and $\beta_1(S)=1$.  Its weighted function is
\begin{equation}\label{eq:example-path}
 C^-_{G,\mm}
   =(1+u+u^2+u^3)e_4+(u+u^2)e_{3,1}.
\end{equation}
For a direct check, the only monomial types are $(2,1,1)$ and
$(1,1,1,1)$.  Enumerating proper color sets gives
\[
 [x_1^2x_2x_3]C^-=u+u^2,\qquad
 [x_1x_2x_3x_4]C^-=1+5u+5u^2+u^3.
\]
Now $e_{3,1}=m_{2,1,1}+4m_{1,1,1,1}$ and
$e_4=m_{1,1,1,1}$, which gives~\eqref{eq:example-path}.
Symmetry is known from Theorem~\ref{thm:A}; alternatively, one may
check each ordered placement of the squared variable directly.
The step-one dual series is
\begin{equation}\label{eq:example-dual}
 \DD^*_{S,1}
   =u e_4+u^2(e_4+e_{3,1})+u^3(e_4+e_{3,1})+u^4e_4.
\end{equation}
The affine series replaces every $e$ by $h$.  Hawkes's original
series in this example are obtained by removing the leading factor
$u$, since $\kappa(S)=1$.

At $\tau=2$, the additional edge between the levels $0$ and $2$
makes the graph complete.  We have
\[
 A_2(S)=5,\qquad
 \beta_2(S)=2+2+1=5.
\]
In the latter expression the three terms come from the equal pair,
the pairs of values $0,1$, and the pair of values $1,2$.
Equation~\eqref{eq:complete} gives
\begin{align}
 C^-_{G,\mm}
   &=[3]_u[4]_u e_4
     =(1+2u+3u^2+3u^3+2u^4+u^5)e_4,\label{eq:example-complete}\\
 \DD^*_{S,2}
   &=u^5(1+2u+3u^2+3u^3+2u^4+u^5)e_4.\label{eq:example-two}
\end{align}
Again the affine series has the same coefficient polynomial with
$h_4$ in place of $e_4$.  This example shows why the step-$\tau$
value graph must include every pair of levels at distance at most
$\tau$, not just consecutive occupied levels.

\subsection{Scope of the conclusions}
The history formula proves nonnegative integrality and the argument
with color reversal proves palindromicity.  Neither assertion alone
implies unimodality or log-concavity of the elementary coefficient
polynomials.  In particular, no answer to the stronger
log-concavity question in~\cite[Question~4.10]{KiemLee2024} is claimed
here.  Likewise, the positivity theorem applies to Hawkes's functions
used in $q,t$-Catalan theory, but it does not by itself assert a new
coefficient formula for a separate Catalan specialization.  The main
new mechanism is the multiplicity target in the common-history
identity, together with the positive finite-support extraction.

\appendix
\section{Correspondence with the cited formalization}\label{app:sources}

The external input is cited at the fixed revision
\href{https://github.com/openai/math/commit/fd4aeeb}{\texttt{fd4aeeb}}
of the OpenAI mathematics repository~\cite{OpenAILean2026}.
The manuscript~\cite{OpenAI2026} supplies the
broader argument; the exact intermediate results used here are
identified below by module and declaration.  All paths are relative
to \nolinkurl{lean/OAI/Combinatorics/Chromatic/}.  The displayed module
names are hyperlinks to that revision, rather than to a mutable
branch.

\subsection{Graph realization and the multiplicity target}
The module
\sourcefile{Walls/GraphTriangularEmbedding.lean}{Walls/GraphTriangularEmbedding.lean}
contains \decl{graph_triangular_embedding}.  It embeds the natural
unit interval graph in a triangular alphabet with the prescribed
pairings.  The passage to the extended lattice and the negative-edge
sign are expressed by \decl{graphVertex_pair} in
\sourcefile{Histories/GraphCommonHistory.lean}{Histories/GraphCommonHistory.lean}.
The vertex-coordinate inclusion is injective; the source uses
\decl{includeVertices_injective} when extracting graph coefficients.
This is the inclusion of the vertex summand into the extended lattice,
so no linear relation among the vertex-coordinate vectors is introduced
by the anchor coordinates.

The same graph-history module contains \decl{squarefree_rootDegree}.
Its domain is an arbitrary injection from a finite set to the
triangular vertices.  Taking the domain to have one element gives
$u_{b(i)}-a_0\in\CC$.  Lemma~\ref{lem:cone} then takes an integral
linear combination with coefficients $m_i$.  The source's theorem
need not be modified or applied to a noninjective map.  This singleton
use is the exact justification for the multiplicity target.

\subsection{Independent elements and the Laurent parameter}
The definitions \decl{IsIndependent}, \decl{setMonomial}, and
\decl{independentElement} are in
\sourcefile{QuantumTorus/IndependentSetElements.lean}{QuantumTorus/IndependentSetElements.lean}.
They give precisely the unnormalized independent-set sum
\eqref{eq:independent}.  In
\sourcefile{Walls/HNTorusFactor.lean}{Walls/HNTorusFactor.lean},
\decl{LaurentRay.vUnit_zpow} identifies the $k$th power of the
parameter unit with the Hahn monomial at exponent $-k$.
Thus the source's parameter unit is our $v^{-1}$, and a torus pairing
$\Omega(a,b)$ contributes $v^{-\Omega(a,b)}$.  This sign convention
is responsible for $v^{A-2\inv}$ rather than $v^{2\inv-A}$ in
Lemma~\ref{lem:kernel}.

\subsection{The common history and its positivity}
The central declaration is \decl{triangular_common_history} in
\sourcefile{Histories/TriangularCommonHistory.lean}{Histories/TriangularCommonHistory.lean}.
Its parameters include a homogeneous degree, a root degree, and an
arbitrary lattice target satisfying \decl{HasRootDegree} after
subtracting the anchor base.  Its conclusion chooses one history
and then quantifies over homogeneous elementary expressions.
The ordered-product specialization in Theorem~\ref{thm:input} uses
\decl{ElementaryExpr.sequence}; the same specialization is visible
in the proof of \decl{graph_common_history} in the graph-history
module.  Nothing in this step imposes a squarefree target.

The histories end at anchored leaves.  Their profile degree is
recorded by \decl{leafProfile_degree}; the dominant-profile property
is stronger than the nonnegativity and total-degree conditions we
need.  In
\sourcefile{Histories/CommonHistoryLeaves.lean}{Histories/CommonHistoryLeaves.lean},
\decl{leafFintype} gives finiteness and \decl{leafWeight_positive}
gives positivity of the product of branch weights for each leaf.
Finally,
\sourcefile{Walls/RootProducts.lean}{Walls/RootProducts.lean}
defines \decl{LaurentPositive.NatCoeff} to mean equality with a
natural number and \decl{LaurentPositive.Positive} coefficientwise
using this predicate.  It is therefore integer positivity, not just
an order condition over the rationals.

\subsection{Notation and scope of the formalized input}
The total coloring degree $N$ in this paper is the homogeneous degree
of an elementary expression.  It is not the number $\ell$ of weighted
vertices, nor the triangular size $D$.  The source's triangular base
at degree $N$ is our $Na_0$, and its root-degree condition is
membership in the integral monoid~\eqref{eq:root-monoid}.  A leaf's
anchor-coordinate profile is our $\mu(\eta)$.  The target used in the
new proof is $R_{\mm}$, whereas the unweighted graph extraction uses
$\sum_i u_{b(i)}$.

These identifications delimit the proof dependency.  The graph
realization and positive common-history statements are imported from
the cited source; the multiplicity coefficient identity, its finite
support argument, and the Dyck applications are proved in this paper.
Only the imported graph-realization and common-history machinery is
used as formalized input.  The new weighted extension and the Dyck
applications are proved here in ordinary mathematical form; no Lean
formalization of those extensions is claimed.

\section{Exact computational checks}\label{app:checks}

Two finite enumerations provide independent checks of the formulas,
including the distinction between ordered monomials and their symmetric
specializations.  All arithmetic is integral, except for exact rational
arithmetic in the change-of-basis calculation.  None of these computations
is used in the proofs.

\subsection{Weighted natural unit interval graphs}
For a fixed $G$, weight vector $\mm$, and composition
$\alpha=(\alpha_1,\ldots,\alpha_r)$ of $N$, the coefficient
$[\xx^\alpha]C^-_{G,\mm}(\xx;u)$ can be computed by processing the colors
in increasing order.  Let $P_a(\mathbf b)$ record the partial colorings
using colors $1,\ldots,a$, with vertex $i$ having received $b_i$ colors.
Here $0\leq b_i\leq m_i$, and $P_0(\mathbf b)$ is $1$ for
$\mathbf b=\mathbf0$ and $0$ otherwise.  For an independent set $I$,
write $\one_I$ for its incidence vector and put
\[
 \Delta_G(\mathbf b,I)
  =\sum_{\substack{i<j\\ij\in E(G)}}\one_{\{i\in I\}}b_j.
\]
The recurrence is
\begin{equation}\label{eq:computational-recurrence}
 P_a(\mathbf b)
  =\sum_{\substack{I\subseteq[\ell]\text{ independent}\\
                   |I|=\alpha_a,\ \one_I\leq\mathbf b}}
       u^{\Delta_G(\mathbf b-\one_I,I)}
       P_{a-1}(\mathbf b-\one_I).
\end{equation}
The exponent counts precisely the new inversions between color $a$ at
a lower-index vertex and an earlier color at an adjacent higher-index
vertex.  Thus the desired coefficient is $P_r(\mm)$, without any
assumption of symmetry.

We enumerated all weakly increasing extensive functions $h$ on $[\ell]$
and all positive weight compositions $(m_1,\ldots,m_\ell)$ with total
weight at most six.  The number of weighted cases at each total weight is
\begin{center}
\begin{tabular}{lrrrrrrr}
\toprule
Total weight $N$ & 1 & 2 & 3 & 4 & 5 & 6 & Total\\
\midrule
Weighted graphs & 1 & 3 & 10 & 36 & 137 & 543 & 730\\
\bottomrule
\end{tabular}
\end{center}
For each degree, the elementary-to-monomial change-of-basis matrix was
computed by counting zero-one matrices.  Exact inversion then gave
candidate elementary coefficients.  These candidates were checked against
every positive composition of that degree, not only against the partition
representatives used to solve the system.  Zero exponents merely skip
colors and do not affect their comparisons.  The tests also checked
coefficientwise nonnegativity, palindromicity with the exact value of $A$,
and the clique blowup identity.

The enumeration verified 19,903 ordered monomial expansions and 6,857
nonzero partition--statistic coefficient positions across the 730 cases.
No discrepancy was found.

\subsection{Direct \texorpdfstring{step-$\tau$}{step-tau} factorizations}
A separate enumeration considered multisets with entries in
$\{0,1,2,3\}$, total size at most five, and
$\tau\in\{0,1,2,3\}$.  For each distinct multiset word, the pair statistic
\eqref{eq:pair-stat} was evaluated directly and compared with the static
decomposition.  All weakly increasing factor labels in four color
variables were then enumerated.  The affine and dual inequalities were
checked separately, and the resulting ordered monomials were weighted
by the directly computed statistic.

These sums were compared with the predicted elementary and homogeneous
expansions.  The homogeneous products were expanded by enumerating weak
compositions for each factor $h_k$; the affine comparison therefore did
not assume the involution identity.  This computation checked 500
multiset--step cases, 5,456 word statistics, and 271,040 candidate
word--factor-label combinations.  The dual and affine comparisons
comprised 39,016 ordered monomial polynomial checks.  No discrepancy was
found, including in the repeated-value grading shifts.

\begin{thebibliography}{Ope26b}

\bibitem[Haw26]{Hawkes2026}
Graham Hawkes, \emph{Dyck symmetric functions and applications to
$q,t$-Catalan polynomials}, preprint, 2026,
\href{https://arxiv.org/abs/2605.13003v1}{arXiv:2605.13003v1}.

\bibitem[Hwa24]{Hwang2024}
Byung-Hak Hwang, \emph{Chromatic quasisymmetric functions and
noncommutative $P$-symmetric functions},
Transactions of the American Mathematical Society
\textbf{377} (2024), no.~4, 2855--2896,
\href{https://doi.org/10.1090/tran/9096}{doi:10.1090/tran/9096};
\href{https://arxiv.org/abs/2208.09857v2}{arXiv:2208.09857v2}.

\bibitem[KL24]{KiemLee2024}
Young-Hoon Kiem and Donggun Lee, \emph{Birational geometry of generalized
Hessenberg varieties and the generalized Shareshian--Wachs conjecture},
Journal of Combinatorial Theory, Series A \textbf{206} (2024),
article 105884,
\href{https://doi.org/10.1016/j.jcta.2024.105884}{doi:10.1016/j.jcta.2024.105884};
\href{https://arxiv.org/abs/2208.12282v3}{arXiv:2208.12282v3}.

\bibitem[Ope26a]{OpenAI2026}
OpenAI, \emph{Elementary positivity of chromatic quasisymmetric
functions}, OpenAI Math Release preprint, September~24, 2026,
\href{https://github.com/openai/math/blob/fd4aeeb/preprints/Elementary-Positivity-of-Chromatic-Quasisymmetric-Functions-September-24-2026/paper.pdf}{OpenAI mathematics repository}.

\bibitem[Ope26b]{OpenAILean2026}
OpenAI, \emph{Formalization accompanying elementary positivity of
chromatic quasisymmetric functions},
\href{https://github.com/openai/math/tree/fd4aeeb/lean/OAI/Combinatorics/Chromatic}{OpenAI mathematics repository, Lean source},
2026, revision
\href{https://github.com/openai/math/commit/fd4aeeb}{\texttt{fd4aeeb}}.

\bibitem[SW16]{ShareshianWachs2016}
John Shareshian and Michelle~L. Wachs, \emph{Chromatic quasisymmetric
functions}, Advances in Mathematics \textbf{295} (2016), 497--551,
\href{https://arxiv.org/abs/1405.4629v2}{arXiv:1405.4629v2}.

\end{thebibliography}

\end{document}
